% TEMPLATE-MILC-v3.tex - plantilla maestra UNICA de las guias MILC v3.
%
% La usan las cuatro familias de guias, cada una con su driver:
%   · Grados 6-11          -> scripts/build-guias-g11.py
%   · Bebras               -> scripts/build-guias-bebras.py
%   · Semillero            -> scripts/build-guias-semillero.py
%   · Territorio interior  -> scripts/build-guias-territorio-interior.py
%
% Antes habia tres copias casi identicas (~400 lineas iguales, 6 distintas):
% cada mejora habia que aplicarla tres veces, y en la practica no se hacia
% (el motor de recursos vivia solo en dos de ellas). Lo que varia entre
% familias esta parametrizado en 6 placeholders que cada driver llena
% (se nombran aqui SIN su sintaxis real: el builder reemplaza por texto plano,
% tambien dentro de los comentarios, y un valor multilinea rompe el preambulo):
%
%   HEADER_IZQ        encabezado izquierdo de pagina
%   HEADER_DER        encabezado derecho de pagina
%   PORTADA_ETIQUETA  rotulo sobre el numero grande (GRADO/MOMENTO/MODULO)
%   PORTADA_SERIE     linea de serie de la portada
%   PORTADA_ID        identificador de la guia en la portada
%   PORTADA_CIRCULO   el circulito de grado (°), o vacio si no aplica
%   PDF_TITULO        titulo en texto plano (sin \textbf ni comillas TeX) para
%                     los metadatos del PDF (/Title); lo produce lib_xelatex.texto_pdf
%
% Composicion (auditoria 2026-09-03): las bandas de estacion piden espacio
% (\Needspace*) y prohiben el salto justo despues (\nopagebreak), asi no quedan
% huerfanas al pie; las cajas largas (softbox, drawbox, infoband, puentebox) son
% `breakable`, asi una caja de media pagina ya no arrastra todo a la siguiente
% dejando la anterior al 40 %. Todo texto sobre mostaza o verde va en milcNegro
% (blanco sobre #E5B400 da 1,93:1 y sobre #58A923 2,95:1; ninguno pasa WCAG AA).
% El resto de claves las documenta content/guias/_SCHEMA.md. El script valida
% que no queden placeholders sin reemplazar antes de escribir el archivo final.

% Etiquetado PDF (latex-lab): /Lang, estructura y texto alternativo de las
% figuras, para lectores de pantalla. Debe ir ANTES de \documentclass.
\DocumentMetadata{lang=es-CO,tagging=on}
\documentclass[11pt,letterpaper]{article}
% headheight=24pt: la cabecera derecha lleva dos lineas (identidad de la guia
% y estacion actual). headsep se acorta en la misma medida para que el bloque
% de texto quede exactamente donde estaba con headheight=12pt.
\usepackage[margin=2.0cm,headheight=24pt,headsep=13pt]{geometry}
\usepackage[table]{xcolor}
\usepackage{fontspec}
\usepackage[spanish]{babel}
\usepackage{tikz}
% Librerías TikZ para los diagramas de las guías (bloque `recursos:`):
%  · babel  → babel en español vuelve activos caracteres como «>», y TikZ los usa
%             en sus opciones (>=stealth, ->). Sin esta librería, cualquier
%             diagrama con flechas revienta con «\language@active@arg> has an
%             extra }». Debe ir ANTES de que se dibuje nada.
%  · positioning        → sintaxis «below left=2pt and 3pt».
%  · decorations.pathreplacing → llaves ({brace}) para señalar rangos.
\usetikzlibrary{babel,positioning,decorations.pathreplacing}
\usepackage{graphicx}
% Circuitos: el trabajo de electrónica se hace hoy con micro:bit y sus sensores
% integrados (MakeCode, sin cableado), así que no se necesitan esquemáticos.
% Si algún día entran montajes con protoboard, basta descomentar esta línea:
% los diagramas .tex/.tikz declarados en `recursos:` ya se incrustan solos.
% \usepackage{circuitikz}
\usepackage[most]{tcolorbox}
\usepackage{tabularx}
\usepackage{array}
\usepackage{fancyhdr}
\usepackage{titlesec}
\usepackage[document]{ragged2e}  % bandera a la derecha en todo el cuerpo
\usepackage{enumitem}
\usepackage{microtype}
\usepackage{etoolbox}
\usepackage{needspace}
% hyperref al final del preambulo: metadatos del PDF (titulo, autor, idioma,
% familia), marcadores por estacion (\pdfbookmark) y enlaces sin recuadro.
\usepackage[hidelinks,
  pdftitle={Alertas con lógica compuesta — la tabla antes que el código},
  pdfauthor={PhD. Álvaro Cárdenas Orozco},
  pdfsubject={Tecnología e Informática · Grado 8 · Guía 18 · MILC v3},
  pdflang={es-CO},
  bookmarksopen=true,bookmarksnumbered=false]{hyperref}

% Helvetica Neue no trae la flecha U+2192, asi que cada «→» escrito literalmente
% en el YAML se compilaba a NADA, en silencio (462 flechas en 61 guias: rutas del
% tipo «paleta → tipografia → lockup» salian rotas). Mapeamos el caracter a su
% version matematica, que si tiene glifo. Asi el autor puede seguir escribiendo
% «→» comodamente en el YAML.
\usepackage{newunicodechar}
\newunicodechar{→}{\ensuremath{\rightarrow}}
% Mismo problema con los simbolos de marca y vinetas: el «✓» de «una palomita ✓»
% salia como cuadrito vacio en el PDF de todas las guias que lo usaban.
\usepackage{pifont}
\newunicodechar{✓}{\ding{51}}
\newunicodechar{✗}{\ding{55}}
\newunicodechar{✔}{\ding{52}}
\newunicodechar{•}{\textbullet}
\newunicodechar{≥}{\ensuremath{\geq}}
\newunicodechar{≤}{\ensuremath{\leq}}

\setmainfont{Helvetica Neue}
\setsansfont{Helvetica Neue}
\setlength{\parindent}{0pt}
\setlength{\parskip}{6pt}
% Sin particion de palabras: en 6.º-7.º una palabra cortada por guion al final
% de linea («Comuni-cacion») frena la lectura mucho mas que una linea corta.
\hyphenpenalty=10000
\exhyphenpenalty=10000
\pretolerance=2000
\tolerance=3000
\setlength{\emergencystretch}{3em}
\linespread{1.10}
\renewcommand{\arraystretch}{1.25}
\setlist{leftmargin=5mm,itemsep=1mm,topsep=1mm}
\newcolumntype{Y}{>{\RaggedRight\arraybackslash}X}

\definecolor{milcMagenta}{HTML}{E6007E}
\definecolor{milcVino}{HTML}{5A0038}
\definecolor{milcTurquesa}{HTML}{0093A5}
\definecolor{milcVerde}{HTML}{58A923}
\definecolor{milcMostaza}{HTML}{E5B400}
\definecolor{milcOcre}{HTML}{B86600}
\definecolor{milcNegro}{HTML}{241816}
\definecolor{milcGris}{HTML}{F7F7F4}
\definecolor{milcLinea}{HTML}{E8E2DD}
\definecolor{milcGradePrimary}{HTML}{FF6600}
\definecolor{milcGradeDark}{HTML}{9A3E00}
\definecolor{milcGradeSoft}{HTML}{FFE4D0}
\definecolor{milcGradeAccent}{HTML}{CFE7FF}
\definecolor{milcGradeText}{HTML}{FFFFFF}
\color{milcNegro}

\pagestyle{fancy}
\fancyhf{}
% Cabecera de dos lineas: identidad de la guia arriba y, a la derecha y en
% gris, la estacion en curso (\markright desde \milcestacion). Antes de la
% primera estacion la segunda linea va vacia; el \strut mantiene la altura
% para que la linea de arriba no baile entre paginas.
\lhead{\small Tecnología e Informática · Grado 8\\[-1pt]{\footnotesize\strut}}
\rhead{\small Guía 18 · MILC v3\\[-1pt]{\footnotesize\strut\color{milcNegro!65}\rightmark}}
\cfoot{\small\thepage}
\renewcommand{\headrulewidth}{0pt}

% Sobre mostaza (#E5B400) y verde (#58A923) el texto va en milcNegro; sobre el
% resto de la paleta, en blanco. Se usa dentro de fonttitle/fontupper, donde un
% \color posterior manda sobre coltitle.
\newcommand{\milctintasobre}[1]{%
  \ifboolexpr{test{\ifstrequal{#1}{milcMostaza}} or test{\ifstrequal{#1}{milcVerde}}}%
    {\color{milcNegro}}{\color{white}}}

\newtcolorbox{titlebox}[1]{
  enhanced,colback=#1,colframe=#1,arc=7mm,boxrule=0pt,
  left=7mm,right=7mm,top=3mm,bottom=3mm,
  before skip=8pt,after skip=8pt,
  fontupper=\sffamily\bfseries\Large\color{white}
}

\newtcolorbox{softbox}[3]{
  enhanced,breakable,pad at break=3mm,
  colback=#2,colframe=#1,borderline west={4pt}{0pt}{#1},
  arc=2mm,boxrule=.4pt,left=4mm,right=4mm,top=3mm,bottom=3mm,
  before skip=6pt,after skip=8pt,title={#3},coltitle=white,
  colbacktitle=#1,fonttitle=\sffamily\bfseries\milctintasobre{#1},fontupper=\RaggedRight,
  attach boxed title to top left={xshift=3mm,yshift=-2mm},
  boxed title style={arc=2mm, boxrule=0pt}
}

\newtcolorbox{drawbox}[2]{
  enhanced,breakable,pad at break=4mm,
  colback=white,colframe=#1,arc=4mm,boxrule=1pt,
  left=5mm,right=5mm,top=4mm,bottom=4mm,before skip=8pt,after skip=8pt,
  title={#2},coltitle=white,colbacktitle=#1,fonttitle=\sffamily\bfseries\milctintasobre{#1},
  fontupper=\RaggedRight,
  attach boxed title to top left={xshift=4mm,yshift=-2mm},
  boxed title style={arc=3mm, boxrule=0pt}
}

\newtcolorbox{infoband}[1]{
  enhanced,breakable,pad at break=2mm,colback=#1!8,colframe=#1!50,arc=2mm,boxrule=.5pt,
  left=4mm,right=4mm,top=2mm,bottom=2mm,before skip=4pt,after skip=4pt,
  fontupper=\small\RaggedRight
}

% Puente narrativo entre secciones (contrato editorial Conéctate).
% Da continuidad: "Acabas de X. Ahora vas a Y porque Z."
% Visualmente: banda izquierda en color del grado, texto en itálica pequeña.
\newtcolorbox{puentebox}{
  enhanced,breakable,pad at break=2mm,colback=milcGradeSoft,colframe=milcGradeDark,
  borderline west={3pt}{0pt}{milcGradeDark},
  arc=0mm,boxrule=0pt,left=5mm,right=4mm,top=2.5mm,bottom=2.5mm,
  before skip=4pt,after skip=8pt,
  fontupper=\itshape\small\color{milcNegro}\RaggedRight
}
% La flecha va en modo matematico a proposito: Helvetica Neue NO tiene el glifo
% U+2192, asi que \textrightarrow se compilaba a NADA («Missing character: There
% is no → in font Helvetica Neue») y el puente salia sin flecha. En modo
% matematico la flecha viene de la fuente matematica y si se dibuja.
\newcommand{\puente}[1]{%
  \begin{puentebox}%
    \textcolor{milcGradeDark}{$\rightarrow$}\hspace{2mm}#1%
  \end{puentebox}%
}

\newcommand{\linead}[1]{\noindent\color{milcLinea}\rule{#1}{0.5pt}\color{milcNegro}}
\newcommand{\respuesta}[1]{\par\vspace{2mm}\foreach \n in {1,...,#1}{\noindent\color{milcLinea}\rule{\linewidth}{0.5pt}\par\vspace{5mm}}\color{milcNegro}}
\newcommand{\checkbox}{\textcolor{milcGradeDark}{\fbox{\phantom{x}}}\hspace{5pt}}

% ─── Listas aireadas para las guías socioemocionales ────────────────────
% Componer las verificaciones como «\checkbox texto\\» deja el texto que salta
% de línea debajo del cuadrito y apelmaza el bloque. Estos entornos alinean la
% continuación y airean los ítems. Los usa Territorio Interior; cualquier
% familia puede hacerlo.
\newlist{checklista}{itemize}{1}
\setlist[checklista]{
  label=\textcolor{milcGradeDark}{\fbox{\phantom{x}}},
  leftmargin=9mm, labelsep=3.2mm, itemsep=2.4mm,
  topsep=2.5mm, parsep=0pt, partopsep=0pt, before=\RaggedRight
}
\newlist{pasoslista}{enumerate}{1}
\setlist[pasoslista]{
  label=\textcolor{milcGradeDark}{\sffamily\bfseries\arabic*.},
  leftmargin=8mm, labelsep=3mm, itemsep=2.8mm,
  topsep=1mm, parsep=0pt, partopsep=0pt, before=\RaggedRight
}

\newcommand{\phasecard}[4]{
\begin{tcolorbox}[enhanced,colback=#3,colframe=#2,arc=3mm,boxrule=.5pt,height=3.05cm,valign=center,halign=center,left=1.5mm,right=1.5mm,top=2mm,bottom=2mm]
{\sffamily\bfseries\fontsize{22}{24}\selectfont\textcolor{#2}{#1}}\par
{\sffamily\bfseries\small #4}
\end{tcolorbox}
}

% ── Piezas del rediseno 2026 (legibilidad 6.º-11.º) ───────────────────
% Banda de estacion: reemplaza el titlebox macizo. Titulo a la izquierda,
% dato practico (tiempo/modalidad) a la derecha, en una sola linea.
% Nunca huerfana: pide 9 lineas libres (o salta de pagina rellenando con
% \vfill) y prohibe el salto justo despues de la banda. Nueve y no seis:
% la caja partible que sigue necesita ~80pt (titulo adosado, relleno y dos
% lineas) para arrancar en la misma pagina; con menos, tcolorbox la manda
% entera a la siguiente y la banda se queda sola. El penalty 10000 va
% en `after`, ANTES del espacio posterior: un `after skip` (aunque sea 0pt)
% es un \vskip tras la caja, y ese glue es un punto de corte legal que un
% \nopagebreak escrito despues de \end{tcolorbox} ya no cierra. Cada banda
% deja marcador PDF y actualiza la cabecera derecha.
\newcounter{milcestacion}
\newcommand{\milcestacion}[2]{%
\Needspace*{9\baselineskip}%
\stepcounter{milcestacion}%
\pdfbookmark[1]{#2}{milcestacion\themilcestacion}%
\markright{#2}%
\begin{tcolorbox}[enhanced,colback=#1,colframe=#1,arc=2mm,boxrule=0pt,
  left=5mm,right=5mm,top=2.6mm,bottom=2.6mm,before skip=11pt,
  after={\par\nopagebreak\vskip7pt}]
{\sffamily\bfseries\fontsize{15}{18}\selectfont\milctintasobre{#1}#2}
\end{tcolorbox}}

% Tarjeta de paso numerada: sustituye las tablas «Pilar / Aplicacion», que
% eran formato de consulta docente, no de estudiante. El numero va en un
% circulo sobre el borde para que la secuencia se lea de un vistazo.
\newcommand{\milcpaso}[3]{%
\Needspace{4\baselineskip}%
\begin{tcolorbox}[enhanced,colback=white,colframe=milcLinea,arc=1.5mm,boxrule=.9pt,
  left=14mm,right=4mm,top=3mm,bottom=3mm,before skip=4pt,after skip=4pt,
  fontupper=\RaggedRight,
  overlay={\node[anchor=north west,circle,fill=#2,text=white,inner sep=0pt,
    minimum size=7.4mm,font=\sffamily\bfseries\small]
    at ([xshift=3.6mm,yshift=-3.2mm]frame.north west) {#1};}]
#3
\end{tcolorbox}}

% Casilla marcable de tamano util con lapiz (la anterior era \fbox{\phantom{x}},
% de alto variable segun el texto que la seguia).
\newcommand{\milcmarca}{\framebox[3.8mm]{\rule{0pt}{3.4mm}}}

% Fila marcable de lista suelta (checks de la estacion 1).
\newcommand{\milccheckrow}[1]{%
\par\vspace{1.8mm}\noindent
\makebox[6.4mm][l]{\raisebox{-.55mm}{\textcolor{milcNegro}{\milcmarca}}}%
\parbox[t]{\dimexpr\linewidth-6.4mm\relax}{\RaggedRight #1}\par}

% Celda marcable dentro de la rubrica (tabla de autoevaluacion). Misma
% sangria francesa, pero pensada para vivir dentro de una columna X.
\newcommand{\milcopcion}[1]{%
\makebox[5.2mm][l]{\raisebox{-.55mm}{\textcolor{milcNegro}{\milcmarca}}}%
\parbox[t]{\dimexpr\linewidth-5.2mm\relax}{\RaggedRight #1}}

% ── Recursos visuales de la guía ──────────────────────────────────────
% Declarados en el bloque `recursos:` del YAML y emitidos por el builder.
% \guiaFigura   → imagen rasterizada o PDF (foto, carta celeste, montaje).
% \guiaDiagrama → fuente TikZ/circuitikz (.tex) incrustada como vector:
%                 es la vía para circuitos y esquemáticos editables.
% [#1] = texto alternativo (alt) · #2 = ruta relativa al .tex · #3 = pie
% El alt llega al PDF etiquetado (/Alt) y lo lee el lector de pantalla.
\newcommand{\guiaFigura}[3][]{%
\begin{center}
\includegraphics[width=0.88\linewidth,height=0.38\textheight,keepaspectratio,alt={#1}]{#2}\\[1.5mm]
{\footnotesize\itshape\textcolor{milcNegro!75}{#3}}
\end{center}
}

\newcommand{\guiaDiagrama}[2]{%
\begin{center}
\input{#1}\\[1.5mm]
{\footnotesize\itshape\textcolor{milcNegro!75}{#2}}
\end{center}
}

\begin{document}
\thispagestyle{empty}

% ── Portada ───────────────────────────────────────────────────────────
% Antes: un rectangulo de color a pagina completa. Se veia bien en pantalla,
% pero en la fotocopiadora del colegio se comia el toner y salia gris sucio.
% Ahora la banda ocupa el tercio superior y el resto va en blanco.
\begin{tikzpicture}[remember picture,overlay]
  \path[fill=milcGradePrimary]
    (current page.north west) rectangle ([yshift=-10.6cm]current page.north east);

  \node[anchor=north west,text=milcGradeText] at ([xshift=2.0cm,yshift=-1.55cm]current page.north west)
    {\sffamily\bfseries\fontsize{12}{14}\selectfont GRADO};
  \node[anchor=north west,text=milcGradeText,opacity=.85] at ([xshift=2.0cm,yshift=-2.35cm]current page.north west)
    {\sffamily\bfseries\fontsize{8.5}{10}\selectfont SERIE GUÍAS · TECNOLOGÍA E INFORMÁTICA};
  \node[anchor=north east,text=milcGradeText,opacity=.85,align=right] at ([xshift=-2.0cm,yshift=-1.55cm]current page.north east)
    {\sffamily\bfseries\fontsize{9}{12}\selectfont I.E. Sor María Juliana\\Cartago, Valle del Cauca};

  \node[anchor=west,text=milcGradeText] (gradonum) at ([xshift=1.85cm,yshift=-7.1cm]current page.north west)
    {\sffamily\bfseries\fontsize{112}{112}\selectfont 8};
  % El circulito de grado (°) solo aplica a las guias de grado («11°»); en
  % Bebras y Semillero el numero grande es un momento/modulo, no un grado.
  % Se posiciona relativo al nodo `gradonum`, no en coordenadas absolutas.
  \draw[milcGradeText,line width=1.6pt]
    ([xshift=2.0mm,yshift=-4.5mm]gradonum.north east) circle (.15cm);

  \node[anchor=west,text=milcGradeText,text width=11.4cm,align=left] at ([xshift=1.5cm,yshift=0.5cm]gradonum.east)
    {
     {\sffamily\bfseries\fontsize{9}{11}\selectfont\MakeUppercase{Período 2 · Lógica, algoritmos y computación física}}\\[3mm]
     {\sffamily\bfseries\fontsize{21}{25}\selectfont\setlength{\baselineskip}{27pt}Alertas\\[4pt]con lógica compuesta}\\[3.5mm]
     {\sffamily\bfseries\fontsize{15}{18}\selectfont Guía 18-8-TIC}
    };
\end{tikzpicture}

\vspace*{9.4cm}
\vfill

{\sffamily\bfseries\fontsize{8}{10}\selectfont\textcolor{milcGradeDark}{LO QUE TE LLEVAS HOY}}

\begin{tcolorbox}[enhanced,colback=white,colframe=milcNegro,arc=2mm,boxrule=1.6pt,
  left=5mm,right=5mm,top=4mm,bottom=4mm,before skip=3pt,after skip=12pt,
  fontupper=\RaggedRight\fontsize{11}{15.5}\selectfont]
Un sistema de alerta en MakeCode que combine dos o tres sensores con Y, O y NO para producir tres niveles: alarma, alerta y aviso. Y una tabla de ocho escenarios, escrita antes de programar, con la salida esperada y la salida real de cada fila, más una nota de los ajustes que hiciste cuando alguna fila no coincidió.
\end{tcolorbox}

\begin{tcolorbox}[enhanced,colback=milcGris,colframe=milcGris,arc=2mm,boxrule=0pt,
  left=5mm,right=5mm,top=3.5mm,bottom=3.5mm,after skip=12pt]
{\sffamily\bfseries\fontsize{8}{10}\selectfont\textcolor{milcNegro!55}{ESTA GUÍA ES DE}}\\[3mm]
\begin{tabularx}{\linewidth}{X p{3.2cm} p{3.2cm}}
\linead{\linewidth} & \linead{3.0cm} & \linead{3.0cm}\\[-1mm]
{\fontsize{8.5}{10}\selectfont\textcolor{milcNegro!55}{Nombre completo}} &
{\fontsize{8.5}{10}\selectfont\textcolor{milcNegro!55}{Grupo}} &
{\fontsize{8.5}{10}\selectfont\textcolor{milcNegro!55}{Fecha}}\\
\end{tabularx}
\end{tcolorbox}

\vfill

\noindent\textcolor{milcLinea}{\rule{\linewidth}{1pt}}\\[2mm]
{\sffamily\bfseries\fontsize{9.5}{12}\selectfont Autor: Dr. Álvaro Cárdenas Orozco}\\[1mm]
{\sffamily\fontsize{8.5}{11}\selectfont\textcolor{milcNegro!55}{Metodología MILC · apertura ancestral · pensamiento computacional · triángulo Dussel–estoicismo–Floridi}}

\newpage

\section*{Apertura: Alertas con lógica compuesta --- la tabla antes que el código}

\begin{softbox}{milcGradeDark}{milcGradeSoft}{Saber ancestral}
En la sesión 6 viste que los mayores nasa de Toribío leen la lluvia en el sapo, las hormigas y la luna. Hay una señal que no se lee sola. Las golondrinas, cuenta don Luis Ardo Ascué, dependen de la hora: si pasan en la mañana, hacia las 8, 9 o 10, es porque va a llover. Si pasan en la tarde, a partir de la una, anuncian otra cosa: problemas sociales, un caso inesperado, un conflicto (Ramos García, Tenorio y Muñoz Yule, 2011). La misma golondrina, con otra hora, es otra alerta. Eso es una condición compuesta: golondrinas \textbf{y} mañana significa lluvia; golondrinas \textbf{y} tarde significa problema. Un sensor solo no decide; decide junto con otro. La \textbf{cara de exclusión}: en Toribío las golondrinas de la tarde han anunciado el asesinato de líderes. El conflicto armado forma parte de lo que ese pueblo lee en el cielo. Los autores advierten que no se trata de probar si aciertan, sino de entender cómo un pueblo ve el mundo. Hoy vas a construir alertas que combinan dos o tres señales. Y una tabla que revise todas las combinaciones antes de dar el sistema por bueno.\par\smallskip{\footnotesize\textcolor{milcNegro!70}{\textbf{Fuente.} Ramos García, C., Tenorio, A. D. y Muñoz Yule, F. (2011). Ciclos naturales, ciclos culturales: percepción y conocimientos tradicionales de los nasas frente al cambio climático en Toribío, Cauca, Colombia. En A. Ulloa (Ed.), \emph{Perspectivas culturales del clima} (pp. 247--274). Universidad Nacional de Colombia.}}
\end{softbox}

\begin{softbox}{milcTurquesa}{milcGris}{Saber contemporáneo}
Una alerta real casi nunca depende de una sola variable. Una alarma de incendio combina humo \textbf{y} temperatura alta, para no sonar cada vez que alguien cocina. Un riego automático abre el agua si la tierra está seca \textbf{y no} está lloviendo. En MakeCode, la categoría Lógica trae tres bloques para eso: «A y B» es verdadero solo si las dos lo son; «A o B», si al menos una; «no A» invierte el valor. Se pueden anidar: «(A y B) o (no C)». Con tres sensores de sí o no hay $2^3 = 8$ combinaciones posibles. Para comprobar el sistema se escribe una \textbf{tabla de escenarios}: una fila por combinación, con la salida que esperas. Después se prueba cada fila en el simulador y se anota la salida real. Si una fila no coincide, hay un error, y se busca con el método de la sesión 7. Un sistema con filas sin probar es un sistema del que no sabes qué hace en esas filas.
\end{softbox}

\begin{softbox}{milcMostaza}{milcGris}{Pregunta-puente}
\textit{Golondrinas y mañana es lluvia; golondrinas y tarde es problema. Si tu alarma de «aula oscura ocupada» se dispara con luz baja, ¿qué segunda señal necesitas para que no suene en el salón vacío?}
\end{softbox}

\begin{softbox}{milcVerde}{milcGris}{Saber y saber hacer}
Al terminar podrás: (1) \textbf{identificar} un problema del colegio que necesite dos o tres señales combinadas; (2) \textbf{crear} la tabla de ocho escenarios con la salida esperada, antes de programar; (3) \textbf{evaluar} las ocho filas en el simulador y corregir las que no coincidan; (4) \textbf{explicar} la diferencia entre una alarma falsa y una alarma muda.
\end{softbox}



\small
\begin{tabularx}{\textwidth}{>{\bfseries}p{3.0cm}Y}
\rowcolor{milcGradePrimary!12}Autor & Dr. Álvaro Cárdenas Orozco\\
DBA & Pensamiento computacional aplicado a sistemas de control con sensores y actuadores (MEN, T\&I)\\
Referentes & Pensamiento computacional · MakeCode / micro:bit · Logos como razón universal\\
Producto final & Un sistema de alerta en MakeCode que combine dos o tres sensores con Y, O y NO para producir tres niveles: alarma, alerta y aviso. Y una tabla de ocho escenarios, escrita antes de programar, con la salida esperada y la salida real de cada fila, más una nota de los ajustes que hiciste cuando alguna fila no coincidió.\\
\end{tabularx}
\normalsize

\puente{En la sesión 7 aprendiste a buscar el error con una hipótesis y una prueba. Hoy combinas los sensores de las sesiones 4 y 6 con Y, O y NO, y pruebas todas las combinaciones. En la sesión 9 vas a dejar el micro:bit midiendo durante tres jornadas.}

\milcestacion{milcGradeDark}{Ruta MILC: así vamos a aprender}
\begin{tabularx}{\textwidth}{>{\centering\arraybackslash}X>{\centering\arraybackslash}X>{\centering\arraybackslash}X>{\centering\arraybackslash}X}
\phasecard{1}{milcVerde}{milcGris}{Escucha\\[-1mm]\small Observo, escucho y pregunto.} &
\phasecard{2}{milcTurquesa}{milcGris}{Entender\\[-1mm]\small Sistematizo y ordeno.} &
\phasecard{3}{milcMagenta}{milcGris}{Hacer\\[-1mm]\small Construyo y aplico.} &
\phasecard{4}{milcVino}{milcGris}{Cierre\\[-1mm]\small Evalúo y reflexiono.}
\end{tabularx}


\puente{Antes de abrir MakeCode, vas a elegir un problema del colegio y decidir qué señales lo detectan. Es la golondrina y la hora: dos señales, no una.}

\milcestacion{milcVerde}{Estación 1 de 3 · Escucha}

\begin{softbox}{milcVerde}{milcGris}{Escena para escuchar}
\textbf{Actividad 1 · IDENTIFICA --- El problema y sus señales} (15 min · individual). (1) Elige un problema del colegio que necesite más de una señal. Ejemplos: aula oscura con gente adentro, luz baja y movimiento; salón que quedó abierto de noche, botón de puerta y hora tardía; sala de sistemas caliente con gente, temperatura alta y movimiento. (2) Escribe qué dos o tres sensores del micro:bit usarías: luz, temperatura, movimiento, botón. (3) Escribe la condición en palabras con Y, O y NO, como «luz baja Y movimiento». (4) Define tres niveles: alarma, alerta y aviso, y qué combinación produce cada uno. (5) Escribe una combinación en la que tu sistema no debería hacer nada.
\end{softbox}

{\sffamily\bfseries\fontsize{8}{10}\selectfont\textcolor{milcNegro!55}{MARCA CUANDO LO HAYAS HECHO}}
\milccheckrow{Elegí un problema que necesita al menos dos señales.}
\milccheckrow{Escribí la condición en palabras con Y, O o NO.}
\milccheckrow{Tengo tres niveles y una combinación sin alerta.}

\begin{infoband}{milcVerde}


\textbf{Lo que va al cuaderno (Actividad 1 · IDENTIFICA).} \textbf{Título:} «Actividad 1 --- El problema y sus señales». \textbf{Formato:} el problema en una línea, los sensores, la condición en palabras con Y, O y NO, y los tres niveles con su combinación. \textbf{Extensión:} media página. \textbf{Sabes que terminaste cuando} puedes decir qué combinación dispara cada nivel y cuál no dispara nada.
\end{infoband}

\puente{Ya tienes el problema y las señales. Ahora aprendes los tres bloques de lógica y, con tu pareja, escribes la tabla de ocho escenarios con lápiz.}

\milcestacion{milcTurquesa}{Estación 2 de 3 · Entender}

\begin{softbox}{milcTurquesa}{milcGris}{Lo que vas a entender}
\textbf{Actividad 2 · CREA --- La tabla de ocho escenarios} (30 min · parejas). (1) Con tu pareja, dibujen una tabla de ocho filas y cinco columnas: sensor 1, sensor 2, sensor 3, salida esperada, salida real. (2) Llenen las tres primeras columnas con todas las combinaciones de V y F: VVV, VVF, VFV, VFF, FVV, FVF, FFV, FFF. (3) Para cada fila, escriban la salida esperada según su lógica: alarma, alerta, aviso o nada. (4) Revisen que ninguna fila quedó sin salida. (5) Marquen las dos filas que les parecen más difíciles de decidir y escriban por qué.
\end{softbox}

\milcpaso{1}{milcTurquesa}{\textbf{Los tres bloques de lógica.} «A y B» solo es verdadero si las dos lo son: luz baja y movimiento. «A o B» es verdadero si al menos una: puerta abierta o ventana abierta. «No A» invierte: no está lloviendo. Se combinan anidando: «(luz baja y movimiento) o (no puerta cerrada)».}
\milcpaso{2}{milcTurquesa}{\textbf{Por qué ocho filas.} Cada sensor de sí o no tiene dos estados. Con dos sensores hay cuatro combinaciones; con tres, ocho; con cuatro, dieciséis. La tabla las escribe todas para que ninguna quede sin decidir.}
\milcpaso{3}{milcTurquesa}{\textbf{La tabla antes que el código.} Si escribes la tabla después de programar, tiendes a confirmar lo que ya hiciste. Si la escribes antes, piensas la lógica sin el programa encima, y el programa tiene que cumplirla.}
\milcpaso{4}{milcTurquesa}{\textbf{Cómo se hace, paso a paso.} (1) Tres columnas de sensores con las ocho combinaciones. (2) Una salida esperada por fila. (3) Ninguna fila vacía. (4) Marcar las difíciles. (5) Programar. (6) Probar y anotar la salida real.}

\begin{drawbox}{milcTurquesa}{La lógica compuesta, en una mirada}
\textbf{A y B:} verdadero solo si las dos. \\[1mm] \textbf{A o B:} verdadero si al menos una. \\[1mm] \textbf{No A:} invierte el valor. \\[1mm] \textbf{Escenario:} una combinación de V y F de los sensores. \\[1mm] \textbf{Tabla:} ocho filas, esperado y real.
\end{drawbox}

\begin{softbox}{milcMostaza}{milcGris}{Errores comunes y su corrección}
(1) \textbf{Alarma falsa}: suena cuando no debe, casi siempre por un O donde iba un Y. (2) \textbf{Alarma muda}: no suena cuando debería, casi siempre por un Y donde iba un O. (3) \textbf{Filas sin probar}: el sistema se declara listo con dos filas de ocho. (4) \textbf{Y y O confundidos}: la condición en palabras dice una cosa y el bloque otra. (5) \textbf{Una sola salida}: tres combinaciones distintas y el mismo ícono para todas.
\end{softbox}

\begin{infoband}{milcTurquesa}


\textbf{Lo que va al cuaderno (Actividad 2 · CREA).} \textbf{Título:} «Actividad 2 --- La tabla de ocho escenarios». \textbf{Formato:} la tabla de ocho filas y cinco columnas con las tres de sensores y la de salida esperada llenas, la de salida real vacía, y las dos filas difíciles marcadas. \textbf{Extensión:} media página. \textbf{Sabes que terminaste cuando} las ocho filas tienen salida esperada y ninguna combinación de V y F se repite.
\end{infoband}

\puente{Ya tienes la tabla con la salida esperada. Toca programar el sistema, probar las ocho filas y corregir las que no coincidan.}

\milcestacion{milcMagenta}{Estación 3 de 3 · Hacer}

\begin{softbox}{milcMagenta}{milcGris}{Cómo vas a construir tu producto}
\textbf{Actividad 3 · EVALÚA --- Programar y probar las ocho filas} (25 min · individual). (1) En MakeCode, dentro de «para siempre», guarda cada lectura en una variable y compárala con su umbral de la sesión 6. (2) Arma un «si… si no, si… si no» con tres ramas usando «y», «o» y «no» según tu tabla, y un ícono distinto por nivel. (3) En el simulador, produce cada una de las ocho combinaciones y anota la salida real en la quinta columna. (4) Si una fila no coincide, aplica el método de la sesión 7: hipótesis, un cambio, prueba. (5) Anota en una nota corta qué fila falló, qué cambiaste y por qué.
\end{softbox}

\milcpaso{1}{milcMagenta}{\textbf{Tu entrega tiene tres piezas.} (1) El programa con tres niveles de alerta, compartido con enlace o archivo .hex. (2) La tabla de ocho escenarios con esperado y real. (3) La nota de ajustes, aunque diga «ninguna fila falló».}
\milcpaso{2}{milcMagenta}{\textbf{Producir una fila en el simulador.} Para «luz baja», tapa el sensor o mueve el control de luz del simulador. Para «movimiento», agita o usa «al agitar». Para el botón, presiónalo. Cada fila se produce con las manos, no se supone.}
\milcpaso{3}{milcMagenta}{\textbf{Falsa y muda.} Si suena en la fila FFF, es alarma falsa: revisa si pusiste O donde iba Y. Si no suena en la fila VVV, es alarma muda: revisa si pusiste Y donde iba O. La tabla te dice cuál de las dos tienes.}
\milcpaso{4}{milcMagenta}{\textbf{Ninguna fila sin probar.} Las dos filas obvias, todo verdadero y todo falso, casi siempre funcionan. Los errores viven en las mixtas: VFV, FVF. Esas son las que la tabla te obliga a mirar.}

\begin{drawbox}{milcMagenta}{Antes de entregar}
\checkbox Dos o tres sensores con sus umbrales de la sesión 6. \\[1mm] \checkbox Lógica con al menos dos de los tres bloques: y, o, no. \\[1mm] \checkbox Tres niveles con un ícono distinto cada uno. \\[1mm] \checkbox Tabla escrita antes de programar, con las ocho filas. \\[1mm] \checkbox Las ocho filas probadas en el simulador con salida real. \\[1mm] \checkbox Nota de ajustes con la fila que falló y el cambio. \\[1mm] \checkbox Programa compartido con enlace o archivo .hex.
\end{drawbox}

\begin{softbox}{milcMagenta}{milcGris}{Plantilla de trabajo}
\textbf{Problema.} \textit{Quiero detectar … cuando …} \\[1mm] \textbf{Señales.} \textit{Sensor 1 …, sensor 2 …, sensor 3 …} \\[1mm] \textbf{Lógica.} \textit{Alarma si … y …; alerta si … o …; aviso si no …} \\[1mm] \textbf{Filas.} \textit{Las ocho, con esperado y real.} \\[1mm] \textbf{Ajuste.} \textit{La fila … falló porque …; cambié …}
\end{softbox}

\begin{infoband}{milcMagenta}


\textbf{Lo que va al cuaderno (Actividad 3 · EVALÚA).} \textbf{Título:} «Actividad 3 --- Programar y probar las ocho filas». \textbf{Formato:} la tabla completa con la columna de salida real llena, y la nota de ajustes con la fila que falló, el cambio y el porqué. \textbf{Extensión:} media página. \textbf{Sabes que terminaste cuando} las ocho filas tienen salida real, esperado y real coinciden, y anotaste qué cambiaste si alguna falló.
\end{infoband}

\puente{Lo que entregas hoy: el programa con tres niveles de alerta, la tabla de ocho escenarios con esperado y real, y la nota de ajustes. La prueba de calidad: las ocho filas tienen salida real anotada y ninguna quedó sin probar.}

\milcestacion{milcMostaza}{Producto de la sesión}
\textbf{Sistema de alerta con tres niveles + tabla de ocho escenarios con esperado y real + nota de ajustes}.
\begin{softbox}{milcMostaza}{milcGris}{Criterios de calidad}
(1) El sistema combina dos o tres sensores con al menos dos de los bloques y, o, no. (2) Produce tres niveles de alerta que se distinguen a simple vista. (3) La tabla tiene las ocho combinaciones y ninguna fila quedó sin salida esperada. (4) Las ocho filas fueron probadas en el simulador y tienen salida real. (5) Esperado y real coinciden, o la nota dice qué se cambió para que coincidan. (6) El programa está compartido con enlace o archivo .hex.
\end{softbox}

\puente{Antes de cerrar, mira lo que hiciste desde cinco lados. Un sistema probado en ocho filas se puede defender; uno probado en las dos filas obvias, no.}

\milcestacion{milcVino}{Evaluación liberadora · cinco dimensiones}

\begin{softbox}{milcVino}{milcGris}{Sentido de la evaluación}
Evaluar no es solo revisar una nota. Es reconocer qué aprendí, qué cambió en mí, qué emociones manejé, cómo me relaciono con la ciudad y la comunidad, y qué le dejaré a quien viene detrás.
\end{softbox}

\textbf{Mi reflexión liberadora en cinco dimensiones.} Completa cada frase iniciada:

\small
\begin{tabularx}{\textwidth}{>{\bfseries\RaggedRight\arraybackslash}p{3.4cm}Y>{\RaggedRight\arraybackslash}p{4.6cm}}
\rowcolor{milcVino}\color{white}Dimensión & \color{white}\bfseries Mi respuesta (frame starter) & \color{white}\bfseries Algo que descubrí\\
1. Personal & Lo que más cambió en mí fue \linead{6.5cm} & \linead{4.2cm}\\
2. Emocional & La emoción más fuerte fue \linead{2.5cm} y la regulé \linead{3cm} & \linead{4.2cm}\\
3. Ciudadana & En democracia esto importa porque \linead{6cm} & \linead{4.2cm}\\
4. Local (Cartago) & En mi barrio sería útil para \linead{6.5cm} & \linead{4.2cm}\\
5. Intergeneracional & A mi abuela/abuelo le diría \linead{6.5cm} & \linead{4.2cm}\\
\end{tabularx}
\normalsize

\begin{infoband}{milcVino}
\textbf{Para la próxima sustentación.} Vuelve a esta tabla en seis meses. Verás que las respuestas cambian — no porque lo aprendido cambie, sino porque tú cambias. La evaluación liberadora es una conversación contigo a través del tiempo.
\end{infoband}


\puente{Para terminar, tres ideas sobre lo que queda fuera de la tabla, contadas con palabras de esta guía: Dussel sobre lo que está más allá del sistema, Séneca sobre los miedos que no aprietan, y la Onlife Initiative sobre la sobrecarga de avisos.}

\milcestacion{milcGradeDark}{Tres ideas para cerrar · Dussel, un estoico y Floridi}

\begin{softbox}{milcVino}{milcGris}{Enrique Dussel · lente del nosotros}
\textit{Toda persona y todo pueblo están siempre más allá del sistema que intenta abarcarlos.}

{\footnotesize\itshape\color{milcNegro!70}--- idea tomada de Enrique Dussel, contada con palabras de esta guía. No es una cita textual.}

Tu tabla de ocho filas es un sistema cerrado. El mundo no: alguien entra al salón con una linterna, y esa combinación no estaba. Probar las ocho filas es lo mínimo; saber que hay una novena es lo honesto.

\textbf{¿Qué situación real no cabe en ninguna fila de mi tabla?}
\end{softbox}
\linead{\linewidth}\\[1mm]
\linead{\linewidth}

\begin{softbox}{milcMostaza}{milcGris}{Estoicismo · Séneca · cuidado interior}
\textit{Son más las cosas que nos asustan que las que de verdad nos aprietan.}

{\footnotesize\itshape\color{milcNegro!70}--- idea tomada de Séneca, contada con palabras de esta guía. No es una cita textual.}

Una alarma falsa asusta sin que pase nada. Un O donde iba un Y hace sonar el sistema por cualquier cosa, y la gente aprende a no hacerle caso. La tabla te muestra en qué fila estás asustando de más.

\textbf{¿Qué alerta me asustó esta semana por algo que al final no apretaba?}
\end{softbox}
\linead{\linewidth}\\[1mm]
\linead{\linewidth}

\begin{softbox}{milcTurquesa}{milcGris}{Luciano Floridi y la Onlife Initiative · lente de la infoesfera}
\textit{La abundancia de información también produce sobrecarga, distracción y olvido.}

{\footnotesize\itshape\color{milcNegro!70}--- idea tomada de Luciano Floridi y la Onlife Initiative, contada con palabras de esta guía. No es una cita textual.}

Un sistema que avisa por todo no informa: satura. Tres niveles bien separados dicen qué tan grave es; un solo ícono para todo obliga a adivinar. Diseñar la lógica es decidir cuándo callar.

\textbf{¿Cuál de mis alertas debería callarse para que las otras se oigan?}
\end{softbox}
\linead{\linewidth}\\[1mm]
\linead{\linewidth}

\begin{infoband}{milcNegro}
\textbf{Nota para el docente.} Las tres ideas están contadas con palabras de esta guía a partir del pensamiento de cada autor; no son citas textuales. De 9.º en adelante se trabajan con la obra y su referencia.
\end{infoband}

\milcestacion{milcVino}{Mi autoevaluación · marca una casilla por fila}

{\small
\setlength{\extrarowheight}{2.2pt}
\begin{tabularx}{\textwidth}{>{\RaggedRight\arraybackslash\bfseries}p{3.0cm}YYY}
\rowcolor{milcVino}\color{white}\bfseries Criterio & \color{white}\bfseries Lo logré & \color{white}\bfseries En proceso & \color{white}\bfseries Necesito apoyo\\[.6mm]
Comprensión del tema & \milcopcion{Explico las ideas con mis palabras.} & \milcopcion{Reconozco las ideas, falta precisión.} & \milcopcion{Necesito repasar lo básico.}\\[.8mm]
\rowcolor{milcGris}Pensamiento computacional & \milcopcion{Usé los pilares al armar mi producto.} & \milcopcion{Usé algunos pilares.} & \milcopcion{Necesito releer la estación 2.}\\[.8mm]
Aplicación y producto & \milcopcion{Mi producto cumple los criterios.} & \milcopcion{Está armado, con ajustes pendientes.} & \milcopcion{Aún está en construcción.}\\[.8mm]
\rowcolor{milcGris}Reflexión 5 dimensiones & \milcopcion{Respondí cada una con honestidad.} & \milcopcion{Respondí algunas con detalle.} & \milcopcion{Mis respuestas son muy generales.}\\[.8mm]
Triángulo de pensamiento & \milcopcion{Conecté las 3 voces con mi trabajo.} & \milcopcion{Conecté al menos una.} & \milcopcion{No las relaciono con mi proceso.}\\[.8mm]
\end{tabularx}}

\vspace{3mm}
\textbf{Hoy entendí que:}\\[1mm]
\linead{\linewidth}\\[1mm]
\linead{\linewidth}

\textbf{Mi compromiso para la próxima sesión es:}\\[1mm]
\linead{\linewidth}\\[1mm]
\linead{\linewidth}

\begin{softbox}{milcMostaza}{milcGris}{Cita de despedida · Marco Aurelio}
\textit{``El obstáculo en el camino se vuelve el camino. Lo que estorba al camino se vuelve camino.''}\\[1mm]
Cada error de hoy, bien observado, se convierte en pieza del camino que pisarás mañana.
\end{softbox}

\small
\begin{tabularx}{\textwidth}{>{\bfseries}p{3.6cm}Y}
\rowcolor{milcGradeDark!12}Nombre completo: & \linead{10cm}\\
Fecha: & \linead{4cm} \quad Curso/grupo: \linead{4cm}\\
Mi firma: & \linead{9cm}\\
\end{tabularx}
\normalsize

\begin{infoband}{milcGradeDark}
\textbf{Para la próxima sesión.} Llega con tu sistema y la tabla de ocho escenarios. En la sesión 9 vas a dejar el micro:bit midiendo una variable del colegio durante tres jornadas, con bitácora. La lógica de hoy decide qué alerta muestra mientras tú no estás mirando.
\end{infoband}

\end{document}
